\SourceLeaf{397}%
die Mittelpunkte von \(AC,BD\) und \(EF\). Bei ganz beliebig gewählter Abszissenachse und beliebigem Ursprung bezeichne ich die Koordinaten dieser neun Punkte mit \(a,a'\quad b,b'\quad c,c'\quad d,d'\) und so weiter. Da nun die drei Punkte \(A,B,E\) auf einer Geraden liegen, erfüllen ihre Koordinaten die folgende Bedingungsgleichung
\[
a(e'-b')+b(a'-e')+e(b'-a')=0
\]
Ebenso gilt, da \(ADF,BCF,DCE\) Geraden sind,
\[
\begin{aligned}
a(f'-d')+d(a'-f')+f(d'-a')&=0\\
b(c'-f')+c(f'-b')+f(b'-c')&=0\\
c(e'-d')+d(c'-e')+e(d'-e')&=0
\end{aligned}
\]

Addiert man diese vier Gleichungen, erhält man
\[
(a+c)(e'+f'-b'-d')+(b+d)(a'+c'-e'-f')+(e+f)(b'+d'-a'-c')=0
\]
oder, da offenbar
\[
\begin{array}{lll}
\tfrac12(a+c)=g,&\tfrac12(b+d)=h,&\tfrac12(e+f)=i\\
\tfrac12(a'+c')=g',&\tfrac12(b'+d')=h',&\tfrac12(e'+f')=i'
\end{array}
\]
gilt,
\[
g(i'-h')+h(g'-i')+i(h'-g')=0
\]
Dies ist die Bedingungsgleichung dafür, dass \(G,H,I\) auf einer Geraden liegen.
\par\bigskip\centerline{\rule{.25\textwidth}{.3pt}}
